\section{实验}
\label{sec:experiments}

\begin{table*}[t]
\caption{\textbf{三个 OLAT 基准全部 18 个场景上的比较}（5,691 张测试图像、相同目标、原始测试标定、不进行测试时拟合；最优用\textbf{粗体}，次优用\underline{下划线}）。迭代次数：RNG 30k+70k，OLAT-GS 30k+30k，GS$^3$ 100k，SSD-GS 100k（SSS 上为 60k），LiSA 38k。}
\label{tab:main}
\centering\small
\setlength{\tabcolsep}{3.8pt}
\begin{tabular}{@{}l ccc ccc ccc ccc@{}}
\toprule
 & \multicolumn{3}{c}{NRHints，真实 (7)} & \multicolumn{3}{c}{GS$^3$，合成 (6)} & \multicolumn{3}{c}{SSS-GS，合成 (5)} & \multicolumn{3}{c}{全部场景 (18)} \\
\cmidrule(lr){2-4}\cmidrule(lr){5-7}\cmidrule(lr){8-10}\cmidrule(l){11-13}
方法 & PSNR$\uparrow$ & SSIM$\uparrow$ & LPIPS$\downarrow$ & PSNR$\uparrow$ & SSIM$\uparrow$ & LPIPS$\downarrow$ & PSNR$\uparrow$ & SSIM$\uparrow$ & LPIPS$\downarrow$ & PSNR$\uparrow$ & SSIM$\uparrow$ & LPIPS$\downarrow$ \\
\midrule
RNG~\cite{fan2025rng} & 20.94 & .828 & .203 & 22.14 & .895 & .132 & 37.10 & .983 & .028 & 25.83 & .893 & .131 \\
OLAT-GS~\cite{kuang2024olat} & 27.26 & \underline{.886} & \underline{.114} & 27.11 & .940 & .079 & \textbf{40.74} & \textbf{.990} & \textbf{.016} & 30.95 & .933 & .075 \\
GS$^3$~\cite{bi2024gs3} & \underline{27.28} & .884 & .116 & \underline{31.97} & \underline{.967} & \underline{.049} & 36.81 & .983 & .028 & 31.49 & \underline{.939} & .069 \\
SSD-GS~\cite{zheng2026ssdgs} & 27.20 & .882 & .117 & 31.53 & .965 & .050 & 37.79 & \underline{.986} & \underline{.023} & \underline{31.59} & \underline{.939} & \underline{.068} \\
\midrule
\textbf{LiSA（本文）} & \textbf{29.28} & \textbf{.909} & \textbf{.100} & \textbf{32.40} & \textbf{.970} & \textbf{.038} & \underline{40.63} & \textbf{.990} & \textbf{.016} & \textbf{33.47} & \textbf{.952} & \textbf{.056} \\
\bottomrule
\end{tabular}
\end{table*}


\subsection{实验设置}
\label{sec:setup}
\paragraph{数据。}
我们使用三个公开 OLAT 基准的所有场景：NRHints~\cite{zeng2023nrhints} 的七个真实拍摄场景（每个场景含 522--1500 张训练图像和 66--200 张测试图像）、\gsthree~\cite{bi2024gs3} 的六个 HDR 合成场景（2000/400），以及 SSS-GS~\cite{dihlmann2024sss} 的五个次表面散射场景（500/500）。
我们遵循各基准的约定（图像分辨率为 512\,px，SSS-GS 为 256\,px；HDR 使用白色背景和 $\gamma=2.2$ 显示），并使用全部官方训练图像及全部 5,691 张测试图像。

\paragraph{基线。}
我们与四种近期公开代码、从点光源拍摄数据逐场景训练的高斯方法比较：\gsthree~\cite{bi2024gs3}、SSD-GS~\cite{zheng2026ssdgs}、OLAT Gaussians（OLAT-GS）~\cite{kuang2024olat} 和 RNG~\cite{fan2025rng}。
各方法均使用其发布的代码和默认训练流程重新训练（60k--100k 次迭代；我们的方法为 38k）。
所有方法均使用前景掩码合成目标图像；OLAT-GS 还将其用于表面阶段，\ours 还将其用于不透明度损失。
\gsthree、SSD-GS 和 OLAT-GS 在真实拍摄数据上优化训练位姿，所有方法均不拟合测试相机或光源。
RNG 和 OLAT-GS 在若干场景（如 Drums、Lego）上低于 21\,dB，可能是因为这些场景不满足其假设（黑色背景的 LDR 数据；代理表面）；我们如实报告这些结果。
\todo{[E5, E6] 获得稳定复现后加入 SSS-GS~\cite{dihlmann2024sss}（我们的实验在五个场景中的三个上发散），并为每种基线复现一个已发表配置。}

\paragraph{指标与协议。}
我们在完整图像上报告 PSNR、SSIM~\cite{wang2004image} 和 VGG-LPIPS~\cite{zhang2018unreasonable}，先对每个场景的测试视图求平均，再对场景求平均，并采用 Wilcoxon 符号秩检验分析逐场景差异。
所有预测都对照同一组 8 位目标图像评分（基线均值的变化最多为 0.005\,dB）；仅在前景掩码内计算 PSNR 时，各基准内部的排名不变（见补充材料）。
主要结果对每种方法、每个场景使用一次训练，消融实验使用三个随机种子。
开发过程中，\ours 的早期版本曾在这些基准上评估（见补充材料）；最终配置依据从训练集中留出的视图选择，所有场景共享该配置。

\subsection{与先进方法比较}
\label{sec:exp_main}
\ours 在 18 个场景上获得了最佳平均 PSNR、SSIM 和 LPIPS（\cref{tab:main}），领先 SSD-GS 1.89\,dB，且相对每种基线的逐场景优势均具有统计显著性（$p\le0.016$）。
没有一种基线在所有基准上都表现一致：OLAT-GS 在 SSS-GS 上最佳，但在 \gsthree 场景上落后 \ours 5.3\,dB；\gsthree 在自身场景上排名第二，却在 SSS-GS 上排名最后。
\ours 是唯一在每个基准上都位列前二的方法，并在全部三个基准上获得最佳 LPIPS。
不过，优势集中于少数场景：三个真实拍摄场景贡献了总优势的 46\%；在六个 \gsthree 场景中，\ours 和 \gsthree 各有三个场景的 PSNR 最佳（\ours 在五个场景的 LPIPS 上最佳）。
在标定精确的 11 个合成场景上，\ours 领先 SSD-GS 1.77\,dB（$p=0.014$），但相对 OLAT-GS 的优势不显著（$p=0.12$）。

最大提升出现在移动投射阴影和穿过半透明物体的光占主导的场景：相较最佳基线，Hotdog 提高 1.90\,dB，Translucent 提高 1.91\,dB。
在 \cref{fig:comparison} 中，花瓶把手和热狗的阴影落在地面及盘子上的大型高斯上；逐图元可见性使阴影变得破碎，而读取图集的像素接收点能重现其轮廓。
\ours 在 Drums（$-1.09$\,dB）和 AnisoMetal（$-0.85$\,dB）上落后，因为其细窄或各向异性高光无法被我们的各向同性瓣函数充分表达。

在真实拍摄数据上，\ours 平均领先 2.0\,dB，但这一优势完全来自 Cat、Fish 和 Pixiu；其余四个场景平均落后 \gsthree 0.3\,dB。
这里，固定标定下的评分不仅取决于外观，也取决于配准：官方位姿存在小幅误差，若优化训练位姿时不锚定重建坐标系，重建可能相对测试相机发生偏移。
我们的规范约束精修（\cref{sec:implementation}）从构造上避免了这一问题。
对每张测试渲染图施加最优二维平移进行对齐（见补充材料），即可看出其影响：基线平均需要 4.0--4.6\,px 的平移，\ours 仅需 1.7\,px；\ours 的优势缩小到 0.7\,dB，在 Cat 和 Pixiu 上仍较大，但在 Fish 上消失。
\torun{E2：不进行位姿精修的 \ours，以及对所有方法仅允许旋转的测试位姿对齐。}

\begin{figure*}[t]
\centering
\setlength{\tabcolsep}{0.8pt}
\renewcommand{\arraystretch}{0.5}
\newcommand{\cw}{0.131\textwidth}
\begin{tabular}{@{}ccccccc@{}}
\small 测试视图 & \small 真值 & \small \ours（本文） & \small SSD-GS & \small \gsthree & \small OLAT-GS & \small RNG \\[2pt]
\includegraphics[width=\cw]{figures/comparison/translucent/full.png} & \includegraphics[width=\cw]{figures/comparison/translucent/GT.png} & \includegraphics[width=\cw]{figures/comparison/translucent/LiSA-staged.png} & \includegraphics[width=\cw]{figures/comparison/translucent/SSD-GS.png} & \includegraphics[width=\cw]{figures/comparison/translucent/GS3.png} & \includegraphics[width=\cw]{figures/comparison/translucent/OLAT-Gaussians.png} & \includegraphics[width=\cw]{figures/comparison/translucent/RNG.png} \\
\footnotesize Translucent & \footnotesize PSNR & \footnotesize \textbf{34.36} & \footnotesize 32.28 & \footnotesize 32.58 & \footnotesize 29.93 & \footnotesize 27.92 \\[3pt]
\includegraphics[width=\cw]{figures/comparison/hotdog/full.png} & \includegraphics[width=\cw]{figures/comparison/hotdog/GT.png} & \includegraphics[width=\cw]{figures/comparison/hotdog/LiSA-staged.png} & \includegraphics[width=\cw]{figures/comparison/hotdog/SSD-GS.png} & \includegraphics[width=\cw]{figures/comparison/hotdog/GS3.png} & \includegraphics[width=\cw]{figures/comparison/hotdog/OLAT-Gaussians.png} & \includegraphics[width=\cw]{figures/comparison/hotdog/RNG.png} \\
\footnotesize Hotdog & \footnotesize PSNR & \footnotesize \textbf{33.97} & \footnotesize 31.93 & \footnotesize 32.08 & \footnotesize 26.85 & \footnotesize 28.49 \\[3pt]
\includegraphics[width=\cw]{figures/comparison/pikachu/full.png} & \includegraphics[width=\cw]{figures/comparison/pikachu/GT.png} & \includegraphics[width=\cw]{figures/comparison/pikachu/LiSA-staged.png} & \includegraphics[width=\cw]{figures/comparison/pikachu/SSD-GS.png} & \includegraphics[width=\cw]{figures/comparison/pikachu/GS3.png} & \includegraphics[width=\cw]{figures/comparison/pikachu/OLAT-Gaussians.png} & \includegraphics[width=\cw]{figures/comparison/pikachu/RNG.png} \\
\footnotesize Pikachu（真实） & \footnotesize PSNR & \footnotesize 30.44 & \footnotesize 29.58 & \footnotesize \textbf{30.71} & \footnotesize 29.46 & \footnotesize 16.31 \\[3pt]
\includegraphics[width=\cw]{figures/comparison/fail_drums/full.png} & \includegraphics[width=\cw]{figures/comparison/fail_drums/GT.png} & \includegraphics[width=\cw]{figures/comparison/fail_drums/LiSA-staged.png} & \includegraphics[width=\cw]{figures/comparison/fail_drums/SSD-GS.png} & \includegraphics[width=\cw]{figures/comparison/fail_drums/GS3.png} & \includegraphics[width=\cw]{figures/comparison/fail_drums/OLAT-Gaussians.png} & \includegraphics[width=\cw]{figures/comparison/fail_drums/RNG.png} \\
\footnotesize Drums & \footnotesize PSNR & \footnotesize 30.90 & \footnotesize 31.32 & \footnotesize \textbf{31.93} & \footnotesize 22.23 & \footnotesize 16.23 \\[3pt]
\end{tabular}
\caption{\textbf{测试视图上的定性比较}（框选区域的局部裁剪；PSNR 对完整视图计算）。各视图均取自对应场景中，\ours 相对最强基线的逐视图优势中位数。各行依次为：两个以投射阴影为主的 \gsthree 场景、一个各方法表现接近的真实拍摄场景，以及我们落后最多的合成场景。SSD-GS 和 \gsthree 的逐图元阴影使花瓶把手及热狗的移动阴影变得破碎；读取图集的像素接收点能够重现它们。\ours 未能重现 Drums 镲片上的细窄光带。更多场景（包括 Fish）见补充材料。}
\label{fig:comparison}
\end{figure*}


\subsection{图集的贡献}
\label{sec:exp_atlas}
\Cref{tab:ablation} 在五个场景组成的测试组上对图集进行消融，覆盖不透明、金属、毛发、半透明和次表面散射材质，其中包含我们落后最多的三个场景；完整模型在不同种子间的波动为 0.11\,dB。
仅使用最细图集层级损失 0.44\,dB，将学习得到的可见性替换为 \cref{eq:moment} 的矩检验则损失 0.24\,dB。
这些单种子差异仅提示两者均有作用；目前尚未在相同接收点下，将图集可见性与逐图元可见性或普通阴影图比较，这将由 E3 完成。

\begin{table}[t]
\caption{\textbf{图集消融}：在五个场景（Drums、AnisoMetal、Translucent、FurScene、Soap）上按场景求平均。三个种子的行报告种子均值和 PSNR 标准差；单种子行与种子 0 的完整模型（31.53\,dB）比较。$^\dagger$容量匹配的 MLP（49.1k 对 49.6k 参数），读取接收点编码、位置和光照/视线方向，但不读取图集；阴影仍来自图集。}
\label{tab:ablation}
\centering\small
\setlength{\tabcolsep}{3.5pt}
\begin{tabular}{@{}lcccc@{}}
\toprule
变体 & 种子数 & PSNR$\uparrow$ & SSIM$\uparrow$ & LPIPS$\downarrow$ \\
\midrule
完整模型 & 3 & 31.51{\tiny$\pm$0.11} & .954 & .047 \\
无光传输项 & 3 & 30.15{\tiny$\pm$0.11} & .948 & .055 \\
光传输 $\rightarrow$ 局部残差$^\dagger$ & 3 & 31.54{\tiny$\pm$0.03} & .954 & .047 \\
单层图集（$K{=}1$） & 1 & 31.09 & .952 & .049 \\
仅矩检验（$g\equiv 0$） & 1 & 31.29 & .952 & .048 \\
逐高斯可见性$^\ddagger$ & 3 & \TBD & \TBD & \TBD \\
硬阴影图$^\ddagger$ & 3 & \TBD & \TBD & \TBD \\
\bottomrule
\end{tabular}
\\[2pt]{\footnotesize$^\ddagger$\torun{E3（相同接收点；光传输保留图集）；单种子行也需三个种子。}}
\end{table}

\begin{table}[t]
\caption{\textbf{光传输项在哪些场景重要。}完整模型相对去除光传输项、相对容量匹配局部残差的逐场景 PSNR 差异（三个种子配对差异的均值$\pm$标准差），以及完整模型分配给光传输项的线性辐亮度比例（8 个测试视图的均值）。}
\label{tab:paired}
\centering\small
\setlength{\tabcolsep}{4pt}
\begin{tabular}{@{}lccc@{}}
\toprule
场景 & 相对无光传输 & 相对局部残差 & 光传输占比 \\
\midrule
Drums & $+0.10${\tiny$\pm$0.02} & $-0.09${\tiny$\pm$0.02} & 6\% \\
FurScene & $+0.02${\tiny$\pm$0.03} & $+0.02${\tiny$\pm$0.02} & 7\% \\
AnisoMetal & $+1.42${\tiny$\pm$0.69} & $-0.18${\tiny$\pm$0.75} & 7\% \\
Translucent & $+2.29${\tiny$\pm$0.23} & $+0.41${\tiny$\pm$0.26} & 17\% \\
Soap & $+2.99${\tiny$\pm$0.30} & $-0.29${\tiny$\pm$0.11} & 58\% \\
\bottomrule
\end{tabular}
\end{table}


在全部三个种子上，去除光传输项使 PSNR 降低 1.37\,dB，LPIPS 从 0.047 升至 0.055；下降幅度随材质而变（\cref{tab:paired}）：Drums 和 FurScene 几乎没有损失，AnisoMetal 损失 1.42\,dB，Translucent 和 Soap 分别损失 2.29 和 2.99\,dB。
前三个场景中，模型将 6--7\% 的线性辐亮度分配给光传输项，Translucent 为 17\%，Soap 为 58\%（在全部 18 个场景中，不透明合成物体为 6--9\%，次表面散射物体为 20--58\%）。
这一分配由学习获得，并非物理分解。

为区分图集本身与额外辐亮度路径的作用，我们用参数容量匹配的 MLP 替代图集光传输项。
该 MLP 读取接收点的编码、位置以及光照和视线方向，但不读取图集；阴影仍由图集提供。
这一局部残差在平均指标上与完整模型相当（31.54 对 31.51\,dB），表明对多数所测材质而言，光传输项主要提供一条不受阴影衰减、非负的辐亮度路径，局部网络同样能够提供它。
图集仅在 Translucent 上领先：光在到达可见表面前明显穿过物体（$+0.41\pm0.26$\,dB，三个种子均为正，但不显著）。
局部残差则在 Soap 上领先（$-0.29\pm0.11$\,dB），在 Drums 上略微领先；在 AnisoMetal 上，完整模型在种子间也更不稳定（26.4--27.9\,dB）。
因此，我们不将总体优势归因于光传输项，而是推测它来自光源空间可见性和优化流程，这将由 E3 与 E4 检验。
相对于光源定义的设计应在光源外推时最为重要，但现有基准并未检验这一点：测试光源与最近训练光源的夹角中位数仅为 $1.1^\circ$。
\torun{E1（\cref{tab:holdout}）：在消融测试组上使用三个种子；需要验证、而不能预设的结论是：图集的性能下降小于局部残差。}

\begin{table}[t]
\caption{\textbf{留出的光照方向}（\torun{E1}）。在官方测试视图和训练集中光照方向落入留出 30$^\circ$ 角度组的视图上，报告五个消融场景的平均 PSNR；各方法均去除这些视图后重新训练。$^\dagger$三个种子。}
\label{tab:holdout}
\centering\small
\begin{tabular}{@{}lccc@{}}
\toprule
方法 & 官方测试 & 留出光源 & 下降 \\
\midrule
\gsthree & 31.25 & \TBD & \TBD \\
SSD-GS & 31.01 & \TBD & \TBD \\
OLAT-GS & 28.67 & \TBD & \TBD \\
\ours，局部残差$^\dagger$ & 31.54 & \TBD & \TBD \\
\ours$^\dagger$ & 31.51 & \TBD & \TBD \\
\bottomrule
\end{tabular}
\end{table}

\begin{figure}[t]
\centering
\setlength{\tabcolsep}{0.8pt}
\renewcommand{\arraystretch}{0.5}
\newcommand{\lw}{0.322\linewidth}
\begin{tabular}{@{}ccc@{}}
\small 显示域损失 & \small 辐射度量分阶段流程 & \small 真值 \\[2pt]
\includegraphics[width=\lw]{figures/loss_domain/lego/control_crop.png} & \includegraphics[width=\lw]{figures/loss_domain/lego/ours_crop.png} & \includegraphics[width=\lw]{figures/loss_domain/lego/gt_crop.png} \\
\footnotesize 23.02\,dB & \footnotesize \textbf{30.25\,dB} & \footnotesize Lego \\[3pt]
\includegraphics[width=\lw]{figures/loss_domain/drums/control_crop.png} & \includegraphics[width=\lw]{figures/loss_domain/drums/ours_crop.png} & \includegraphics[width=\lw]{figures/loss_domain/drums/gt_crop.png} \\
\footnotesize 24.10\,dB & \footnotesize \textbf{30.10\,dB} & \footnotesize Drums \\[3pt]
\end{tabular}
\caption{\textbf{损失计算域与几何。}两次训练共享初始化、训练流程、随机状态及所有其他设置，仅损失计算域不同：显示编码辐亮度，或退火预热后的线性辐亮度。训练 16k 次迭代，PSNR 在 200 个留出训练视图上计算。两个场景的全部 200 个视图均得到改善。}
\label{fig:loss_domain}
\end{figure}


\subsection{优化流程的贡献}
\label{sec:exp_optimization}
\paragraph{损失计算域。}
在初次训练失败的 Lego 和 Drums 上，两次各 16k 迭代的训练共享初始化、预热阶段、训练流程和随机状态，仅损失不同：对照组始终使用显示编码辐亮度，我们的方法遵循阶段 I 和 II。
在从训练集中留出的 200 个视图上，显示域训练得到的几何发生坍塌（\cref{fig:loss_domain}），而按辐射度量域分阶段的流程使 Lego 的 PSNR 从 23.02 升至 30.25\,dB，Drums 从 24.10 升至 30.10\,dB（LPIPS 分别为 0.122$\to$0.057 和 0.060$\to$0.034）。
每个视图均得到改善，两场景的提升中位数分别为 7.2 和 6.0\,dB。
该研究仅涵盖两个 HDR 场景、一个种子和 16k 次迭代；在真实和 SSS-GS 场景上，\ours 较早的显示域版本高出 0.2--0.4\,dB，但该版本还存在其他差异（见补充材料）。
\torun{E4：在消融测试组上使用三个种子，比较 (a) 始终使用显示域损失，(b) 无预热，(c) 阶段 II 使用像素接收点，(d) 无阶段 III，(e) 在 $\Gamma$ 内加入偏移量。}

\paragraph{接收点与精修。}
在 Lego、Drums 和 Soap 的冻结 30k 几何上（补充材料表~S6），以显示域而非线性辐亮度域进行精修，可改善两种接收点的全部三项指标（高斯的 PSNR 提高 $+0.20$\,dB，像素提高 $+0.28$\,dB）。
像素接收点的 SSIM 和 LPIPS 最佳（0.976 和 0.029），高斯接收点的 PSNR 则高出 0.19\,dB；出于感知质量的考虑，我们采用像素接收点。

\subsection{计算开销}
\label{sec:exp_cost}
在其余任务均空闲的 GPU 上对 Drums、FurScene 和 Soap 进行性能测量（补充材料表~S5）时，\ours 的训练速度为 \gsthree 和 SSD-GS 的 1.6--5.5 倍；Drums 上使用 228k 个高斯，而两种基线均约为 550k。
全部 18 个场景在较不统一的计时条件下，平均训练时间为 21 分钟，基线为 55--107 分钟。
即使每帧更换光源，计入光源渲染通道后，整体渲染仍可达到 63--104 帧每秒，保持交互速度；不过，它也使 \ours 在小场景上比 \gsthree 和 SSD-GS 更慢。

\subsection{局限性}
\label{sec:limitations}
各向同性瓣函数无法充分表达拉丝金属和细窄光带：在 Drums 上，真值中被标记为窄高光的像素仅占 0.27\%，却贡献了我们相对 \gsthree 多出的平方误差的 72\%。
矩检验对混合覆盖区域仅提供较弱先验，中心深度也会错误定位大型倾斜高斯；自然的改进方向包括切比雪夫检验或四矩检验~\cite{donnelly2006variance,peters2015moment}，以及使用交点深度。
单一透视图集假设光源位于物体外部，光传输项忽略立体角因子和入射点距离衰减，而仅经由光源不可见区域到达接收点的光，超出了单次光源空间汇集所能表达的范围。
